\documentclass[11pt]{article}
\usepackage[a4paper,margin=25mm]{geometry}
\usepackage{amsmath,amssymb,amsthm,mathtools,booktabs,hyperref,xurl}
\setlength{\emergencystretch}{2em}
\newtheorem{theorem}{Theorem}
\newtheorem{lemma}[theorem]{Lemma}
\newtheorem{proposition}[theorem]{Proposition}
\newtheorem{corollary}[theorem]{Corollary}
\newtheorem{remark}[theorem]{Remark}
\newcommand{\Z}{\mathbb Z}\newcommand{\Q}{\mathbb Q}
\title{Signed binomial witnesses force their own positivity\\and save five quadratic coordinates}
\author{\(\mathbb Z\mathbb Q\) Diophantine Project}
\date{23 November 2026}
\begin{document}
\maketitle

\begin{abstract}
In the corrected Pell certificate for central-binomial divisibility,
the two auxiliary inputs \(w,s\) need no positivity wrappers once
\(s\ne0\) is enforced. We prove this for every integer modulus
\(M\ge16\), uniformly in the prescribed index. The proof classifies
all signs of both Pell parameters. The negative first-sequence branch
turns the power-of-two congruence into its multiplicative inverse;
an explicit size bound excludes that branch and every negative input.
Replacing two three-coordinate positivity wrappers by the single
quadratic guard \(s s_{\rm inv}=1\) saves five coordinates and one
residual. This local result is unconditional. Its application to
undecidability retains the corrected source's cited coding theorem and
every arithmetic hypothesis of the rational integer interface.
\end{abstract}

\section{The exact signed statement}

For an integer \(a\) with \(|a|\ge2\), define the polynomial Pell
sequences by
\[
 X_0(a)=1,\quad X_1(a)=a,\qquad
 Y_0(a)=0,\quad Y_1(a)=1,\qquad
 Z_{j+2}(a)=2aZ_{j+1}(a)-Z_j(a).
\]
They satisfy
\[
 X_j(a)^2-(a^2-1)Y_j(a)^2=1,
 \quad X_j(-a)=(-1)^jX_j(a),
 \quad Y_j(-a)=(-1)^{j-1}Y_j(a).
\]
We use the elementary Pell classification, index congruences, and
signed trace lemma proved in \cite{Pell}; the extension to negative
parameters needed here is proved explicitly below.

\begin{theorem}[Automatic positivity]\label{main}
Let \(r,M,b\in\Z\) satisfy
\[
 r\ge1,\qquad M\ge16,\qquad 2\le b\le2M.
\]
Let \(w,s,c,d,k,h,\gamma,j,o\) be integers, let
\(i\in\Z\setminus\{0\}\), and let \(\tau,F,\zeta\in\Q\).
Assume \(s\ne0\), and put
\begin{equation}\label{definitions}
 Y=sM,\quad B=2r+1,\quad
 a=(wM+1)rY,\quad p=2wr^2Y^2M.
\end{equation}
Suppose
\begin{align}
 k&=r+1+h(p-1),&\tau^2&=(p^2-1)k^2+1,\label{small}\\
 k^2-4(c-kY)^2&>0,\label{close}\\
 d&=bw+(a-2)c+(4a-5)\gamma,\label{linear}\\
 d^2&=(a^2-1)c^2+1,\label{norm1}\\
 F^2&=(a^2-1)i^2c^4+1,\label{norm2}\\
 (d+oF)^2&=(G^2-1)(B+jc)^2+1,
 &G&=a+F^2(c+1-a),\label{norm3}\\
 \zeta^2&=(a-1)(d-1).\label{root}
\end{align}
Then
\[
 w>0,\qquad s>0,\qquad c>0,\qquad bw=2^B.
\]
Consequently the signed block has exactly the same projected
\((r,M,b)\)-solutions as the positive binomial certificate of
\cite{Binomial}. In particular it forces
\(M\mid\binom{2r}{r}\).
\end{theorem}

The final consequence uses the previously proved signed reconstruction:
once \(w,s>0\), the retained equations restore \(c>B\),
\(h>0\), \(d>0\), and \(\gamma>0\); positive \(i,F,j,o\)
can be chosen anew. No new integer predicate is implicit in this
statement. The reciprocal coordinate used in the full rational tables
is discussed in Section~\ref{implementation}.

\section{The first Pell index for either sign of the parameter}

\begin{lemma}\label{signedindex}
Let \(a,c,d,B\) be integers with \(|a|\ge2\), \(B\ge3\)
odd, and \(|c|>B\). Assume \eqref{norm1}--\eqref{norm3} with
\(i\in\Z\setminus\{0\}\), \(j,o\in\Z\), and \(F\in\Q\).
Put \(A=|a|\). Then
\[
 |c|=Y_B(A),\qquad |d|=X_B(A).
\]
If \eqref{root} also holds, then \(d=X_B(a)\).
\end{lemma}
\begin{proof}
Write \(C=|c|\). By Pell classification at \(A\), the first norm
gives \(C=Y_t(A)\), \(|d|=X_t(A)\) for some \(t\ge2\).
Thus \(C\ge2A\) and \(t<C\). Since a rational number with
integer square is an integer, the second norm makes \(F\) an integer.
It is nonzero, and classification gives
\[
 |i|C^2=Y_q(A),\qquad |F|=X_q(A).
\]
The square-divisibility step-down of \cite{Pell} yields \(C\mid q\),
so \(q\ge C>t\).

The bridge has absolute value at least two for every sign of \(a,c\).
If \(c=C>0\), then \(c+1-a\ge C+1-A>0\), and \(F^2\ge1\)
implies \(G\ge c+1\ge5\). If \(c=-C<0\), then
\(c+1-a\le1-A<0\), so \(G\le c+1\le-3\).
The norm congruences also give
\[
 G\equiv1\pmod C,\qquad G\equiv a\pmod {|F|}.
\]
Since \(0<B<C\), the integer \(B+jc\) is nonzero. Signed Pell
classification at \(G\) therefore writes the last norm as
\[
 B+jc=\varepsilon Y_v(G),\qquad
 d+oF=\varepsilon'X_v(G)
\]
with \(v\ge1\) and independent signs \(\varepsilon,\varepsilon'\).
Reduction modulo \(C\) gives \(v\equiv\pm B\pmod C\).
Reduction modulo \(|F|\), followed by
\(X_v(a)=\operatorname{sgn}(a)^vX_v(A)\), gives
\[
 X_v(A)\equiv\pm X_t(A)\pmod{X_q(A)}.
\]
The signed trace lemma yields \(v\equiv\pm t\pmod{2q}\), hence
\(B\equiv\pm t\pmod C\). Because \(0<B,t<C\), either
\(B=t\) or \(B=C-t\). The second alternative is impossible:
\(C\equiv t\pmod2\) makes \(C-t\) even, while \(B\) is odd.
Thus \(t=B\).

Finally \(|d|=X_B(A)>1\). If \(a>0\), equation \eqref{root}
forces \(d\ge1\); if \(a<0\), it forces \(d\le1\).
Therefore \(d=\operatorname{sgn}(a)X_B(A)=X_B(a)\), since \(B\)
is odd.
\end{proof}

\section{Both small-parameter signs and a uniform size bound}

\begin{proof}[Proof of Theorem~\ref{main}]
The inequality \eqref{close} first forces \(k\ne0\) and gives
\begin{equation}\label{half}
 \left|\frac ck-Y\right|<\frac12.
\end{equation}
Since \(|Y|=|s|M\ge16\), it follows that \(c\ne0\).

We first exclude \(w=0\). Then \(p=0\), and the first norm makes
\(\tau\) an integer and gives \(k=\pm1\), \(\tau=0\).
The integer \(c/k-Y\) has absolute value less than \(1/2\),
so \(c=kY\). Now \(a=rY\), and the first large norm would give
\[
 d^2=(rY^2)^2-Y^2+1.
\]
This lies strictly between \((rY^2-1)^2\) and \((rY^2)^2\):
the two positive gaps are \((2r-1)Y^2\) and \(Y^2-1\).
That is impossible for an integer square. Hence \(w\ne0\).

Write
\[
 \sigma=\operatorname{sgn}(w),\quad
 \epsilon=\operatorname{sgn}(s),\quad
 U=|w|M\ge16,\quad Z=|s|M\ge16,
\]
and put
\begin{equation}\label{AP}
 A=|a|=rZ(U+\sigma),\qquad
 P=|p|=2r^2Z^2U.
\end{equation}
In particular \(A\ge2\), \(P\ge8192r^2\), and
\begin{equation}\label{rough}
 P-1>5r+1,\qquad 2A<P^2,\qquad P>A.
\end{equation}
For the middle inequality use \(A\le rZ(U+1)\) and
\(P^2/(2rZ(U+1))=2r^3Z^3U^2/(U+1)>1\).
The others follow directly from \(U,Z\ge16\).

The congruence for \(k\) has modulus
\(|p-1|=P-\sigma>2(r+1)\), so \(|k|\ge r+1\).
Equation \eqref{half} gives
\[
 |c|>|k|(Z-1/2)\ge(r+1)(Z-1/2)>B.
\]
Lemma~\ref{signedindex} now proves, with \(C=|c|\),
\begin{equation}\label{largevalues}
 C=Y_{2r+1}(A),\qquad d=X_{2r+1}(a).
\end{equation}

Again \(\tau\) is an integer. Pell classification of its norm gives
\(|k|=Y_\ell(P)\) for an integer \(\ell\ge1\).
The elementary recurrence bounds and \eqref{rough} give
\[
 C<(2A)^{2r}<P^{4r}.
\]
Since \(|k|<C\), the bound
\(Y_\ell(P)>P^{\ell-1}\) for \(\ell\ge2\) excludes
\(\ell\ge4r+1\). Thus \(\ell\le4r\).

If \(\sigma=1\), then
\(Y_\ell(P)\equiv\ell\pmod{P-1}\). If \(\sigma=-1\),
then
\(Y_\ell(P)\equiv(-1)^{\ell-1}\ell\pmod{P+1}\).
These are the recurrence specialized at \(1\) and \(-1\).
Because \(\ell\le4r\) and the modulus is greater than \(5r+1\),
the congruence \(k\equiv r+1\pmod{p-1}\) forces
\begin{equation}\label{smallvalue}
 \ell=r+1,\qquad k=\sigma^r K,
 \qquad K=Y_{r+1}(P)>0.
\end{equation}
Equation \eqref{half} then yields
\begin{equation}\label{signc}
 \operatorname{sgn}(c)=\epsilon\sigma^r,
 \qquad \left|\frac CK-Z\right|<\frac12.
\end{equation}

We need a uniform estimate which works for \(U-1\) as well as \(U+1\).
For real \(x\ge2\), write
\[
 f_j(x)=\frac{Y_j(x)}{(2x)^{j-1}}.
\]
The normalized Pell lemma in \cite{Binomial} gives
\[
 1-\frac{j-2}{4x^2}\le f_j(x)\le1\qquad(j\ge2).
\]
It also proves that \(f_j(x)\) is nondecreasing in \(x\) and
strictly decreasing in \(j\) from \(j=2\) onward. Since \(P>A\),
\(f_{2r+1}(A)<f_{r+1}(P)\). The leading terms in their ratio give
\[
 \rho=\left(\frac{(U+\sigma)^2}{U}\right)^r,
 \qquad \frac CK=\rho\frac{f_{2r+1}(A)}{f_{r+1}(P)}.
\]
Moreover
\[
 f_{2r+1}(A)>1-\frac{1}{2Z^2(U-1)^2}.
\]
Indeed \((2r-1)/(4r^2)<1/2\) and \(A\ge rZ(U-1)\).
Using \(f_{r+1}(P)\le1\), we obtain the convenient loose bound
\begin{equation}\label{ratio}
 0<\rho-\frac CK<\frac{\rho}{Z^2(U-1)^2}.
\end{equation}
The denominator is at least \(16\), and
\(\rho>(U-2)^r\). Equations \eqref{signc} and \eqref{ratio}
therefore imply
\[
 Z>\frac{15}{16}(U-2)^r-\frac12
   >\frac78(U-2)^r.
\]
The final inequality follows from \((U-2)^r\ge14\).
Consequently
\begin{equation}\label{size}
 A>\frac78 r(U-1)(U-2)^r>2U4^r.
\end{equation}
For the last step, divide by \(U4^r\) and use
\[
 \frac78r\frac{U-1}{U}\left(\frac{U-2}{4}\right)^r
 \ge\frac78\cdot\frac{15}{16}\cdot\frac{14}{4}
 =\frac{735}{256}>2.
\]

\section*{The two congruence branches}

For the signed integer parameter \(a\), define
\[
 R_j^-=X_j(a)-(a-2)Y_j(a),\qquad
 R_j^+=X_j(a)+(a-2)Y_j(a).
\]
The polynomial addition and norm identities give
\begin{equation}\label{congruences}
 R_{j+1}^-=2R_j^-+(4a-5)Y_j(a),\qquad R_0^-=1,
 \qquad R_j^-R_j^+=1+(4a-5)Y_j(a)^2.
\end{equation}
Thus, for \(m=4a-5\),
\[
 R_B^-\equiv2^B\pmod m,\qquad
 2^B R_B^+\equiv1\pmod m.
\]
These identities require no sign assumption on \(a\).
As \(B\) is odd, \(Y_B(a)=Y_B(A)=C\), while
\eqref{largevalues} already gives \(d=X_B(a)\).

The input bound implies \(|bw|\le2U\). From \eqref{size},
\begin{equation}\label{modulusbound}
 |m|\ge4A-5>8U4^r-5>4U4^r+1.
\end{equation}
If \(c=-C\), the linear row identifies \(bw\) with
\(R_B^+\) modulo \(m\). Hence
\(bw2^B\equiv1\pmod m\). But
\[
 0<|bw2^B-1|\le4U4^r+1<|m|.
\]
The first inequality uses that the nonzero integer \(bw\) has
absolute value at least two. This is impossible for a multiple of
\(m\). Therefore \(c=C>0\).

The linear row now gives \(bw\equiv2^B\pmod m\). Also
\[
 |bw-2^B|\le2U+2\cdot4^r<4U4^r+1<|m|.
\]
It follows that \(bw=2^B\), so \(w>0\). Equation \eqref{signc}
then forces \(s>0\). The previously proved positive binomial
certificate and signed reconstruction give the final divisibility
statement, as described after Theorem~\ref{main}.
\end{proof}

\section{A necessary nonzero guard}\label{implementation}

The hypothesis \(s\ne0\) cannot simply be omitted. For example,
take \(r=1\), \(M=256\), \(b=16\) and
\[
 w=1,\ s=0,\ a=p=c=0,\ k=h=i=F=d=1,\
 \tau=\zeta=j=o=0,\ \gamma=3.
\]
Then \(Y=0\), \(B=3\), \(G=1\). Every displayed equation
and the strict inequality hold, but
\(256\nmid\binom21\). This is an exact counterexample to deleting
both positivity wrappers with no replacement guard.

In the complete rational tables, the input predicates already make
\(s,w,c,h,\gamma,j,o\) integers; all other integer coordinates are
reconstructed in the established dependency order. Add one rational
coordinate \(s_{\rm inv}\) and the quadratic equation
\begin{equation}\label{guard}
                         s s_{\rm inv}=1.
\end{equation}
This ensures \(|sM|\ge16\) before the Pell argument begins.
The inequality first gives \(k\ne0\), and if \(c=0\) it would
give \((1-4s^2M^2)k^2>0\), a contradiction. Thus \(c\ne0\)
is established independently of any Pell classification. The existing
row \(vE=c^2\) now gives \(v\ne0\), legitimizing the retained
reciprocal predicate and the reconstruction \(i=1/v\in\Z\setminus\{0\}\).
The norm rows then give the required integrality of \(F\) and
\(\tau\). Theorem~\ref{main} applies without a circular sign or
nonzero assumption.

In the corrected Jones source, \(q=b^{5^{60}}\), \(n=q^{16}\),
and \(M=n^2\), while the retained source prefix gives \(b\ge2\).
Thus \(M\ge16\) and \(b\le2M\) hold with room to spare.

\begin{corollary}[Five fewer coordinates]\label{saving}
In each corrected complete quadratic table, delete the rows
\texttt{jones1980.positive.w} and \texttt{jones1980.positive.s}
and their six private coordinates. Add \eqref{guard} and its one
new coordinate. Every other row remains literally unchanged after
coordinate renumbering. This preserves rational solvability, saves
five coordinates and one row, and retains degree four after taking
the sum of the squared residuals.
\end{corollary}
\begin{proof}
An old rational zero has positive integral \(s\), so it extends by
\(s_{\rm inv}=1/s\). Conversely, a new rational zero satisfies the
noncircular reconstruction just proved; Theorem~\ref{main} yields
positive integers \(w,s\). Each has the existing three-triangular-number
positivity representation, so the six old private coordinates can be
restored without changing any retained coordinate. The occurrence
audit verifies that they occur nowhere else. The count is
\(6-1=5\) coordinates and \(2-1=1\) rows.
\end{proof}

Applied after the source correction and the ceiling power circuit,
and after fixed contact offsets for the ordinary-\(abc\) interface,
the exact complete counts are
\begin{center}
\begin{tabular}{lrrrr}
\toprule
Interface & Old variables & New variables & Old rows & New rows\\
\midrule
General &1074&1069&924&923\\
Specified &1254&1249&1080&1079\\
Ordinary rational \(abc\)&786&781&685&684\\
\bottomrule
\end{tabular}
\end{center}
These are conditional undecidability transfers under exactly the
retained rational integer-interface premises and the corrected source's
explicitly cited Jones coding theorem. The sign theorem and the
five-coordinate equivalence themselves require no such external
arithmetic premise. The new result applies to the corrected \(n^6\)
block; no claim is made for the historical \(n^2\) parameter.

\section{Exact artifacts}

The complete tables and their hashes are built and independently
checked by
\begin{verbatim}
python undecidability/source_repair_v93/build_signed_positivity.py
python undecidability/source_repair_v93/verify_signed_positivity.py
\end{verbatim}
The checker audits every discarded coordinate occurrence, every
retained residual, the nonzero guard, counts, support, degree, and
deterministic reproduction. It also checks the polynomial congruence
identities and the zero-\(s\) counterexample exactly. Its finite Pell
examples illustrate the sign cases; the quantified proof is the
argument above.

\begin{thebibliography}{9}
\bibitem{Pell} \(\mathbb Z\mathbb Q\) Diophantine Project,
\emph{An exact three-Pell index certificate and nine fewer quadratic
coordinates}, 23 November 2026.\par
\nolinkurl{undecidability/papers/pell_index_optimization.tex}.
\bibitem{Binomial} \(\mathbb Z\mathbb Q\) Diophantine Project,
\emph{An exact Pell certificate for a power of two and central-binomial
divisibility}, 23 November 2026.\par
\nolinkurl{undecidability/papers/binomial_pell_certificate.tex}.
\bibitem{Corrected} \(\mathbb Z\mathbb Q\) Diophantine Project,
\emph{A corrected quartic source from Jones's binomial theorem},
23 November 2026.\par
\nolinkurl{undecidability/source_repair_v93/corrected_jones_source.tex}.
\end{thebibliography}
\end{document}
